\section{H.264 media profile and SRT transport}\label{sec:video}
\subsection{Codec, container and transport are separate choices}
\req{VID-01}{Protocol 1.0 SHALL carry H.264/AVC video in MPEG-TS over SRT live
mode. HTTP advertises and allocates streams; WebSocket carries status. Video bytes
SHALL not share the telemetry WebSocket. SRT is the proposed transport binding;
H.264 was the explicitly accepted codec choice.}

The video node is the SRT \emph{listener and sender}; Launch Video is the SRT
\emph{caller and receiver}. Caller/listener describes who establishes the
connection, not who produces the video. Connections are unicast on the LAN.
No Internet rendezvous service, browser or external media server is needed.
SRT's receive buffer trades delay for recovery of late/retransmitted packets;
its latency setting is not the complete camera-to-screen latency \cite{srtlatency}.

\begin{longtable}{L{39mm}L{122mm}}
\toprule Item & Required baseline\\\midrule
\endfirsthead
\toprule Item & Required baseline\\\midrule
\endhead
H.264 decoder profile & Constrained Baseline, level at most 3.1, progressive 8-bit 4:2:0. No B-frames. Other profiles require an explicitly advertised extension.\\
Random access & IDR at least once per second; repeat SPS/PPS at every IDR. New subscribers begin with parameter sets and a decodable IDR.\\
MPEG-TS & One program with one H.264 video elementary stream (stream type 0x1B); Annex-B NAL representation, PTS and PCR. Audio absent in 1.0. Receivers use PAT/PMT rather than hard-coded PIDs.\\
Tables and timing & PAT/PMT at startup and at least every 100 ms; PCR interval at most 40 ms. PTS uses the 90 kHz timebase; handle its 33-bit wrap. These are selected application-profile limits.\\
SRT mode & Live, caller/listener, timestamp-based packet delivery, too-late packet drop enabled. Baseline receiver latency 120 ms; the session reports actual negotiated value.\\
Payload & Up to 1316 bytes (seven 188-byte TS packets) per SRT application send; shorter final groups are permitted. Do not append a custom JSON header to media.\\
Connect/read behavior & Connect timeout 3 s; no decoded frame for 2 s is stale, for 5 s unavailable. An explicit closed connection is immediately unavailable. Reconnect uses the session procedure below.\\
Resource boundaries & One capture per device; one encoder per active channel/rendition; fan out encoded data to distinct per-viewer SRT sessions. A slow viewer cannot block other viewers.\\
\bottomrule
\end{longtable}

MPEG-TS parameters and SRT option mappings must be checked against the packaged
media library. FFmpeg documents these muxer and protocol options
\cite{ffmpegformats,ffmpegprotocols}; using another compliant library is allowed.
In particular, an adapter must convert this document's millisecond values to the
units required by that API. FFmpeg SRT latency options use microseconds.

\subsection{Renditions and manifests}
\begin{longtable}{L{27mm}L{43mm}L{39mm}L{47mm}}
\toprule Rendition & Image/FPS ceiling & Proposed video bitrate & Intended use\\\midrule
\endfirsthead
\toprule Rendition & Image/FPS ceiling & Proposed video bitrate & Intended use\\\midrule
\endhead
\code{preview} & Fit within $320\times180$; 10 fps & 250 kbit/s target; 300 kbit/s cap & Eight small remote views\\
\code{main} & Fit within $1280\times720$; 30 fps & 3 Mbit/s target; 4 Mbit/s cap & Selected large channel\\
\bottomrule
\end{longtable}

Preserve aspect ratio with even dimensions; do not distort a 4:3 analog capture.
Do not invent frames to claim a source FPS it does not deliver. Encoder output
and source FPS SHALL both be reported. These are network rendition ceilings,
not assumptions about the capture device's native resolution. A profile may
deliver lower quality when declared in its actual encoding parameters. Full-HD,
audio and additional codecs are extensions, not baseline requirements.

\code{/streams} returns \code{streams}, an array of entries with:
\code{stream_id}, \code{channel_id}, \code{station}, \code{vehicle},
\code{radio_type}, \code{source_mode}, \code{capture_state},
\code{rf_lock}, \code{media_epoch}, and \code{renditions}. A stream ID is
\code{station.channel_id}, for example \code{gs4.booster.analog}; globally
identify a publisher with node/boot as well. Each rendition describes name,
codec, container, transport, actual resolution/FPS, target/cap bitrate and
\code{available} plus reason. A configured disconnected camera remains in the
manifest with unavailable renditions. No hardware camera identifier is needed
by remote clients.

\code{media_epoch} is a UUID for the channel's current capture/encoding timeline.
It changes on camera replacement, capture restart, or encoder/timestamp reset.
Starting a previously unused rendition alone may join the existing timeline;
any change that invalidates its timing requires a new channel epoch and restarts
all sessions for that channel. The service boot and media epoch are independent.
Send \code{streams.changed} before issuing new sessions for a changed epoch.

\subsection{Session allocation, connection and lease}
The client discovers metadata first; it does not guess a media port. To receive
one rendition it POSTs the following body to \code{/video/sessions}:
\begin{lstlisting}
{
  "request_id": "<UUID>", "client_node_id": "<UUID>",
  "client_boot_id": "<UUID>", "mission_id": "<UUID>",
  "expected_server_boot_id": "<UUID>",
  "stream_id": "gs1.sustainer.digital",
  "rendition": "main", "expected_media_epoch": "<UUID>"
}
\end{lstlisting}
\code{request_id} makes session creation idempotent: within a service boot,
repeating the same ID and identical fields returns the original session result.
Different fields return 409. Keep that record until at least 60 seconds after
the session ends; an expired recorded result returns 410, never a fresh session.
Clients SHALL retry an uncertain POST only within this interval with the same ID;
afterward perform cleanup/reconciliation and create a new explicit request.
For the owning client, \code{/streams} MAY include its active session IDs to aid
reconciliation; this optional convenience is not needed for baseline operation.

The service validates mission, boot, channel assignment, source mode, media epoch,
capacity and camera state; starts/reuses the encoder; binds one available UDP
listener in its configured pool; then returns HTTP 201. Do not return success
before a listener is ready. The response includes \code{session_id},
\code{session_token}, \code{server_boot_id}, \code{media_epoch},
\code{stream_id}, \code{rendition}, \code{transport},
\code{lease_ms=30000} and actual encoding parameters. \code{transport} contains
resolvable host, UDP port, \code{mode=caller} (for the client),
\code{srt_streamid}, \code{latency_ms=120}, \code{container=mpegts} and
\code{codec=h264}. Use the selected reachable interface, never \code{0.0.0.0}
as the advertised destination.

The SRT Stream ID is the ASCII string \code{rt1:<session_id>:<session_token>}.
The token is 128 bits of cryptographically random data encoded as 32 lower-case
hex digits. The caller sends it at connection setup; the listener compares it
exactly before accepting one receiver. Stream ID is connection metadata available
to the listener \cite{srtoptions}. This token prevents accidental session mixups;
on this plaintext LAN profile it does not provide encryption or strong peer trust.
Do not include it in logs, DNS records or screenshots.

Each allocated listener serves exactly one client session; additional connection
attempts are rejected. Sharing one source therefore means sharing encoded bytes
across independent SRT sockets, not assuming a single socket automatically
broadcasts. Baseline capacity is four active sessions per video node, with at
most two main and two preview sessions. Unavailable capacity returns 429 and a
one-second retry hint. These limits support the current video wall, including
two promoted channels at the same Away station.

The client renews every 10 seconds with PATCH, body
\code{\{"lease_ms":30000\}}, and header
\code{Authorization: Bearer <session_token>}. Only 30000 is accepted in 1.0.
Expiry uses the server's monotonic clock, extended to 30 seconds from successful
renewal. PATCH success returns the remaining lease; no cross-computer comparison
of absolute expiry timestamps is required. Expiry or DELETE closes that session,
returns its UDP port and stops unused encoding after a two-second grace period.
DELETE uses the same header and returns 204 even after prior deletion; unknown
tokens cannot affect another live session. Invalid tokens return 403.

On media loss, DELETE/release the previous session and allocate a new one with
a new request ID after rechecking boot/epoch. Do not reuse an accepted SRT session
token for a second socket. A server boot change invalidates all leases and tokens.
Loss of the status WebSocket alone need not destroy healthy media immediately;
mark status stale and retain media only while its own lease/frames remain valid.

\begin{figure}[H]\centering
\begin{tikzpicture}[node distance=7mm]
\node[box,text width=30mm] (usb) {USB capture\\one device owner};
\node[box,text width=38mm,right=7mm of usb] (enc) {H.264 encoders\\preview + main};
\node[box,text width=38mm,right=7mm of enc] (srt) {MPEG-TS / SRT\\one socket per viewer};
\draw[flow] (usb)--(enc); \draw[flow] (enc)--(srt);
\node[accentbox,text width=118mm,below=12mm of enc] (api)
{HTTP: allocate / renew / release session\\WebSocket: camera state, manifest, timing and errors};
\draw[optional] (api.north east)--(srt.south);
\node[box,text width=118mm,below=8mm of api] (wall)
{Launch Video: decode selected sources\\main feed can also supply its own thumbnail};
\draw[flow] (srt.east)--++(7mm,0)|-(wall.east);
\draw[flow,<->] (api)--(wall);
\end{tikzpicture}
\caption{Video pixels and control/status travel separately. An HTTP session
establishes a specific channel, rendition and media epoch before SRT begins.}
\end{figure}


\subsection{Promotion, bandwidth and failure recovery}
\req{VID-02}{Launch Video SHALL normally request previews for all eight Away
feeds. Promoting a feed requests its main rendition, retains the preview until
the first main IDR is decoded, then uses that same main decode for its thumbnail
and releases the now-redundant preview session. Demotion performs the reverse.}
\req{VID-03}{No main rendition SHALL be requested for a logical channel that
the publishing station does not possess. A source swap SHALL update the label
and epoch before displaying new pixels. A failed swap SHALL leave the previous
view intact with an explicit error; it SHALL not mislabel old pixels.}

A conservative budget for eight previews at the cap and three main feeds is
$8(0.3)+3(4)=14.4$ Mbit/s of encoded video. Adding a provisional 30\% allowance
for encapsulation/retransmission gives 18.72 Mbit/s. This deliberately counts
previews also shown from main and so overestimates steady-state use. It does not
include Wi-Fi airtime contention, recording transfer, other users or unusually
high loss. It is not a claim of LTU throughput. Measure every Away path and the
shared Launch link under concurrent load; provision at least twice this planned
load as measured sustained application capacity for the initial deployment.

If capacity falls, reduce preview rate/bitrate first, then offer reduced main
quality explicitly. Prioritize serial ingestion, event publication and heartbeat
processing over video work. Separate sockets alone cannot guarantee QoS on a
saturated radio; encoder caps, admission limits and network configuration are
part of acceptance testing. Initial implementation MAY use fixed profiles with
operator quality selection; automatic congestion adaptation is not assumed to
be supplied by SRT.

\subsection{Media timing and recording}
The receiver SHALL use decoder PTS for playback and its own monotonic arrival
clock for frame-staleness display. Capture time means receipt of a USB frame by
the publishing host, not the rocket camera's exposure time. The receiver/RF path
may add unknown latency before USB capture. SRT timestamps do not provide rocket
camera timestamps.

For correlation, each active rendition's status SHALL report \code{media_epoch},
\code{pts_90k} as a decimal unwrapped counter, \code{capture_mono_ms},
\code{capture_time_utc}, and clock uncertainty for one identified encoded frame
at least once per second. TS PTS is that counter modulo $2^{33}$. Two rendition
encoders SHALL derive times from the same capture timeline; frames may be omitted
for preview. Resetting the timeline changes the epoch. This mapping supports
approximate telemetry/video correlation; accurate cross-station overlays remain
unqualified until clock and capture delays are measured.

Local recording SHALL be independent of remote viewer quality. When the operator
enables recording, keep local video/telemetry with source IDs, mission, epochs,
timestamps, clock quality, decoder/codec settings and gap events. Launch recording
can archive what it receives but SHALL not claim to contain frames lost upstream.
Disk-full or recording failure SHALL be visible without blocking live acquisition.
After a video gap, resume on an IDR and reset the decoder as needed; do not display
corrupted inter-frame predictions as a current reliable image.
